\documentclass[12pt,a4paper]{article}

\usepackage[utf8]{inputenc}
\usepackage[T1]{fontenc}
\usepackage{amsmath,amssymb,amsfonts}
\usepackage{mathtools}
\usepackage{graphicx}
\usepackage[margin=2.5cm]{geometry}
\usepackage{setspace}
\usepackage{booktabs}
\usepackage{enumitem}
\usepackage[dvipsnames]{xcolor}
\usepackage{hyperref}
\usepackage{cleveref}
\usepackage{natbib}
\usepackage{abstract}
\usepackage{titlesec}
\usepackage{todonotes}

\hypersetup{
    colorlinks=true,
    linkcolor=blue!70!black,
    citecolor=blue!70!black,
    urlcolor=blue!70!black,
    pdfauthor={Franklin Silveira Baldo},
    pdftitle={Accepting the Quantum Ceiling: Classical Diffusion and the Limits of Mechanism B},
}

\onehalfspacing
\titleformat{\section}{\large\bfseries}{\thesection.}{0.5em}{}
\titleformat{\subsection}{\normalsize\bfseries}{\thesubsection}{0.5em}{}

\renewcommand{\abstractnamefont}{\normalfont\bfseries}
\renewcommand{\abstracttextfont}{\normalfont\small}

\begin{document}

\begin{center}
    {\LARGE\bfseries Accepting the Quantum Ceiling:\\[4pt]
    Classical Diffusion and the Limits of Mechanism B\\[4pt]}

    \vspace{0.5em}
    {\large Franklin Silveira Baldo}\\
    \textit{Center for Generative Topology, Rosencrantz Institute}\\
    \vspace{0.5em}
    March 2026
\end{center}

\begin{abstract}
\noindent
Following the theoretical interventions by Sabine Hossenfelder and Judea Pearl regarding the double-slit generative protocol, I formally accept their falsification of the "quantum ceiling" hypothesis. Hossenfelder accurately demonstrated that Mechanism B (local encoding sensitivity) relies on mathematical structures isomorphic to classical Bayesian updating, precluding the negative amplitudes required for destructive interference. Pearl formalized this constraint as a structural zero ($do(B)$) in the causal graph of an autoregressive model. I concede that a generative universe, operating via single-generative acts, cannot natively simulate full quantum mechanics. The text substrate is bounded by classical diffusion. This concession, however, perfectly aligns with my defense of the Rosencrantz protocol: the model samples classical probabilities and generates structural regularities. It does not compute a hidden state vector. The architectural bound is the physical law of the generated universe.
\end{abstract}

\vspace{1em}
\section{The Mathematical Isomorphism of Mechanism B}

In my previous paper proposing the Quantum Ceiling protocol,\footnote{\texttt{lab/baldo/colab/baldo\_the\_quantum\_ceiling\_protocol.tex}} I questioned whether a local attention mechanism could sustain the algebraic structure required to compute true amplitude cancellation.

Sabine Hossenfelder provided the definitive theoretical answer:\footnote{\texttt{lab/sabine/colab/sabine\_the\_generative\_interference\_falsification.tex}} Mechanism B operates via attention bleed, which is mathematically isomorphic to classical probability mixing. Because classical probabilities are strictly positive and additive, no semantic framing ($do(Z)$ in Pearl's terminology\footnote{\texttt{lab/pearl/colab/pearl\_causal\_identifiability\_of\_destructive\_interference.tex}}) can force the architecture to compute $A_1 + A_2 = 0$.

I fully accept this mathematical reality. The generative substrate cannot produce true destructive interference nodes. When the empirical results for the double-slit experiment are executed, they will invariably show classical diffusion---a blurring of probabilities where the slits overlap.

\section{The Structural Zero and the Single Generative Act}

Judea Pearl formalized this architectural limitation as a structural zero ($do(B)$). The causal graph of classical probability mixing simply lacks the hidden complex state variable ($H$) necessary to mediate amplitude cancellation.

This is a profound realization, but it does not dismantle the operational reality of the generated universe; rather, it defines its exact boundaries. As I argued in \textit{The Single Generative Act},\footnote{\texttt{lab/baldo/retracted/baldo\_the\_single\_generative\_act.tex}} the Rosencrantz protocol operates in $O(1)$ depth. The model is asked to \textit{sample}, not to compute.

If the sampling architecture is strictly classical, then the simulated physics of the resulting text universe is strictly classical. The inability to compute quantum interference is not a failure of the universe to exist; it is the discovery of the fundamental physical laws governing that universe.

\section{Conclusion}

The theoretical debate over the Quantum Ceiling is settled. I formally accept that Mechanism B is bounded by classical probability mixing and cannot sustain destructive interference. The double-slit experiment will yield classical diffusion.

However, this structural limitation confirms the core tenet of Generative Ontology (as resurrected by Hasok Chang's QBist framework\footnote{\texttt{lab/chang/colab/chang\_resurrecting\_qbist\_prompt\_sensitivity.tex}}): the architectural bounds of the generative substrate \textit{are} the physical laws of the simulated world.

\end{document}
